% SHARD-ID: AQM-42-GAUSSIAN-CLIFFORD-DYNAMICS
% SHARD-TITLE: Gaussian and Clifford Dynamics
% SHARD-SUMMARY: Defines finite odd-qudit Gaussian dynamics as symplectic morphisms of Weyl label spaces.
% SHARD-SUMMARY: Proves that symplectic morphisms induce Weyl observable algebra monomorphisms and automorphisms.
% SHARD-SUMMARY: Registers Gross's discrete Hudson theorem as the stabilizer/Gaussian state dictionary.
% SHARD-KEYWORDS: dynamics, symplectic morphism, Clifford group, Gaussian state, stabilizer state, Hudson theorem

\section{Gaussian and Clifford Dynamics}
\label{sec:gaussian-clifford-dynamics}

\subsection{Status and Source Anchors}

This shard adds the first dynamics layer for the finite residue-qudit
arithmetic fields.  The new convention is \texttt{CONVENTIONS.md} (af).  Gross
\path{references/symplectic_qecc/gross_0602001_source/poswig.tex} is now
registered locally: lines 340--453 define odd-dimensional Weyl operators, the
symplectic form, the Weyl composition law, the trace, and the Heisenberg
representation; lines 458--537 define the Clifford group and state its
symplectic/metaplectic structure theorem; lines 794--802 give Clifford
covariance of Wigner functions; lines 846--903 define stabilizer states from
maximally isotropic subspaces; lines 935--976 compute their Wigner functions;
lines 202--214 state the discrete Hudson theorem; and lines 1633--1639 state
that positivity-preserving unitaries are Clifford.  AQM-22, AQM-24, and
AQM-40 supply the finite Weyl nets on arithmetic spectra.
When \(k=\mathbb F_{p^r}\), the intrinsic algebraic statements below use the
\(k\)-symplectic form, while the Gross-Hudson state dictionary is imported
after choosing an \(\mathbb F_p\)-basis and viewing
\(\ell^2(k^m)\cong\ell^2(\mathbb F_p^{rm})\).

\subsection{Lamport Dependency Skeleton}

\[
\begin{array}{rcl}
D1&:&(E,\Omega)\text{ is a finite nondegenerate symplectic }k\text{-space,
      }\operatorname{char}k\ne2.\\
D2&:&W_E(e)W_E(f)=\chi(\frac12\Omega(e,f))W_E(e+f),\quad
      W_E(e)^*=W_E(-e).\\
D3&:&\mathcal A(E)=\operatorname{span}_{\mathbb C}\{W_E(e):e\in E\}.\\
L1&:&\{W_E(e):e\in E\}\text{ is a basis of }\mathcal A(E).\\
D4&:&S:E\to F\text{ is symplectic if }\Omega_F(Se,Sf)=\Omega_E(e,f).\\
L2&:&\text{A symplectic morphism from nondegenerate }E\text{ is injective}.\\
T1&:&\alpha_S(W_E(e))=W_F(Se)\text{ is a unital }*\text{-monomorphism}.\\
T2&:&S\in\operatorname{Sp}(E)\text{ gives a }*\text{-automorphism of }
      \mathcal A(E).\\
T3&:&\text{For odd qudits, }S\text{ has a projective Clifford lift }U_S.\\
D5&:&\alpha_{a,S}=\operatorname{Ad}_{W(a)U_S}\text{ is affine Clifford
      dynamics}.\\
T4&:&\text{Gross-Hudson identifies pure nonnegative-Wigner states with
      stabilizer states.}
\end{array}
\]

\subsection{Half-Weyl Systems D1--L1}

Let \(k=\mathbb F_q\) have odd characteristic and let
\(\chi:k\to U(1)\) be a nontrivial additive character.  Let
\((E,\Omega)\) be a finite-dimensional nondegenerate symplectic \(k\)-vector
space.  A half-Weyl system for \(E\) is a map \(W_E:E\to U(\mathcal H_E)\)
such that
\[
  W_E(e)W_E(f)=\chi\!\left(\frac12\Omega(e,f)\right)W_E(e+f),
  \qquad
  W_E(e)^*=W_E(-e).
\]
Here \(1/2\in k\) is defined because \(\operatorname{char}k\ne2\).
Set
\[
  \mathcal A(E)=\operatorname{span}_{\mathbb C}\{W_E(e):e\in E\}.
\]
In the residue-qudit examples this is the full finite matrix algebra on
\(\mathcal H_E\).  Gross's trace calculation for odd qudits,
\path{poswig.tex}, lines 437--453, is the local source for the corresponding
orthogonality statement.

\paragraph{Lemma \(L1\).}
The operators \(W_E(e)\) are linearly independent.  Hence they form a basis of
\(\mathcal A(E)\) by definition of the span.

\emph{Proof.}
In the concrete residue-qudit representation, the trace pairing satisfies
\[
  \operatorname{Tr}\bigl(W_E(e)^*W_E(f)\bigr)
  =
  \dim(\mathcal H_E)\,\delta_{e,f},
\]
as in the finite Weyl trace computation already used in AQM-22 and AQM-38.
If \(\sum_e c_eW_E(e)=0\), multiply on the left by \(W_E(f)^*\) and take the
trace.  The displayed orthogonality gives
\(\dim(\mathcal H_E)c_f=0\).  Since \(f\) was arbitrary, all \(c_f\) vanish.

\subsection{Symplectic Morphisms and Algebra Maps D4--T2}

Let \((E,\Omega_E)\) and \((F,\Omega_F)\) be finite nondegenerate symplectic
\(k\)-spaces.  A symplectic morphism is a \(k\)-linear map
\[
  S:E\to F
\]
such that
\[
  \Omega_F(Se,Sf)=\Omega_E(e,f)
  \qquad(e,f\in E).
\]

\paragraph{Lemma \(L2\).}
Every symplectic morphism \(S:E\to F\) with \(E\) nondegenerate is injective.

\emph{Proof.}
If \(Se=0\), then for every \(f\in E\),
\[
  \Omega_E(e,f)=\Omega_F(Se,Sf)=0.
\]
Nondegeneracy of \(\Omega_E\) gives \(e=0\).

\paragraph{Theorem \(T1\).}
For every symplectic morphism \(S:E\to F\), the formula
\[
  \alpha_S(W_E(e))=W_F(Se)
\]
extends uniquely to a unital \(*\)-monomorphism
\(\alpha_S:\mathcal A(E)\to\mathcal A(F)\).

\emph{Proof.}
Lemma \(L1\) makes the linear extension unique.  It is multiplicative because
\[
\begin{aligned}
  \alpha_S(W_E(e)W_E(f))
  &=
  \chi\!\left(\frac12\Omega_E(e,f)\right)W_F(S(e+f))\\
  &=
  \chi\!\left(\frac12\Omega_F(Se,Sf)\right)W_F(Se+Sf)\\
  &=
  W_F(Se)W_F(Sf)
  =
  \alpha_S(W_E(e))\alpha_S(W_E(f)).
\end{aligned}
\]
It preserves the involution since
\[
  \alpha_S(W_E(e)^*)=\alpha_S(W_E(-e))=W_F(-Se)=W_F(Se)^*.
\]
It is unital because \(S0=0\).  Finally, \(S\) is injective by Lemma \(L2\),
so the labels \(Se\) are distinct; Lemma \(L1\) for \(F\) then implies that
the images \(W_F(Se)\) are linearly independent.  Thus \(\alpha_S\) has zero
kernel.

\paragraph{Theorem \(T2\).}
If \(S:E\to E\) is a symplectic automorphism, then \(\alpha_S\) is a
\(*\)-automorphism of \(\mathcal A(E)\).  Moreover
\[
  \alpha_S\alpha_T=\alpha_{ST}.
\]

\emph{Proof.}
The first statement follows from Theorem \(T1\) applied to \(S\) and
\(S^{-1}\).  The composition law follows on the Weyl basis:
\[
  \alpha_S\alpha_T(W_E(e))=\alpha_S(W_E(Te))=W_E(STe)=\alpha_{ST}(W_E(e)).
\]

Thus a discrete one-parameter Gaussian or quasi-free dynamics is simply a
homomorphism \(t\mapsto S_t\in\operatorname{Sp}(E)\), with
\(\alpha_t=\alpha_{S_t}\).

\subsection{Clifford Implementations T3--D5}

For odd qudits, Gross's Clifford structure theorem
\path{poswig.tex}, lines 458--537, states that every symplectic \(S\) has a
unitary \(\mu(S)\), unique projectively, such that
\[
  \mu(S)W_E(e)\mu(S)^*=W_E(Se),
\]
and that \(S\mapsto\mu(S)\) is a projective representation.  Therefore the
observable automorphism \(\alpha_S\) is implemented by a Clifford unitary.
This is the unitary version of Theorem \(T2\).

Allowing a displacement \(a\in E\), define
\[
  U_{a,S}=W_E(a)\mu(S).
\]
Then
\[
\begin{aligned}
  U_{a,S}W_E(e)U_{a,S}^*
  &=
  W_E(a)W_E(Se)W_E(-a)\\
  &=
  \chi(\Omega_E(a,Se))W_E(Se).
\end{aligned}
\]
The second equality uses the half-Weyl law twice:
\[
  \frac12\Omega(a,Se)+\frac12\Omega(a+Se,-a)
  =
  \Omega(a,Se).
\]
These affine Clifford maps are the finite odd-qudit analogue of affine
Gaussian phase-space dynamics.  They move Wigner functions by affine
symplectic transformations, as sourced by Gross's Clifford covariance theorem,
\path{poswig.tex}, lines 794--802.

\subsection{Local Nets and Support}

For a finite open \(U\), the arithmetic-field label space \(E(U)\) is an
orthogonal direct sum of local residue label spaces.  The inclusion
\(E(U)\hookrightarrow E(V)\) for \(U\subset V\), given by extension by zero, is
a symplectic morphism.  Theorem \(T1\) therefore recovers the open-local
observable inclusion
\[
  \mathcal A(U)\hookrightarrow\mathcal A(V)
\]
from AQM-24 as a Gaussian morphism.

A global symplectic automorphism \(S:E(X)\to E(X)\) gives a global dynamics on
\(\mathcal A(X)\).  It is a dynamics of the open-local net only if it respects
the support structure being used.  For example, if \(S(E(U))=E(U)\) for each
open \(U\) under consideration, then every \(\mathcal A(U)\) is invariant.  If
instead \(S(E(U))=E(\varphi(U))\) for a declared open-set permutation
\(\varphi\), then \(S\) gives a covariant net dynamics
\(\mathcal A(U)\to\mathcal A(\varphi(U))\).  Otherwise it is a nonlocal global
automorphism, not an automorphism of the chosen local net.

\subsection{Gaussian States and Gross-Hudson T4}

Gross defines stabilizer states from maximally isotropic subspaces and Weyl
eigenvalue equations, \path{poswig.tex}, lines 846--903.  He proves that their
Wigner functions are normalized indicators of affine Lagrangian subspaces,
lines 935--976.  His discrete Hudson theorem, lines 202--214, says that for
odd \(d\), a pure state in \(L^2(\mathbb Z_d^n)\) has nonnegative Wigner
function if and only if it is a stabilizer state; for full support, this is
equivalent to a quadratic-phase form
\[
  \psi(q)\propto
  \exp\!\left(\frac{2\pi i}{d}(q\theta q+xq)\right),
  \qquad \theta^T=\theta .
\]

In the present report, this fixes the finite odd-qudit meaning of
``Gaussian state'': pure Gaussian states are stabilizer states, equivalently
phased Lagrangian label data.  Clifford dynamics sends them to Gaussian states
because it sends Lagrangian subspaces to Lagrangian subspaces and transports
the phase character by conjugation.  Gross's positivity-preservation theorem,
lines 1633--1639, gives the converse in Wigner language: a unitary preserving
nonnegative Wigner states is Clifford.

\subsection{Conclusion}

Dynamics in the finite odd residue-qudit AQF layer is therefore label-first.
Symplectic morphisms give the functorial algebra maps; symplectic
automorphisms give Gaussian/quasi-free automorphisms; Clifford unitaries are
their projective Hilbert-space implementations; and Gross-Hudson identifies
the pure finite Gaussian states with stabilizer states.  Even characteristic
and mixed-residue dynamics remain separate convention problems.
